\documentclass[12pt,a4paper]{report}

\usepackage[utf8]{inputenc}
\usepackage[T1]{fontenc}
\usepackage{geometry}
\usepackage{graphicx}
\usepackage{listings}
\usepackage{xcolor}
\usepackage{hyperref}
\usepackage{fancyhdr}
\usepackage{titlesec}
\usepackage{booktabs}
\usepackage{array}
\usepackage{enumitem}
\usepackage{float}
\usepackage{caption}
\usepackage{subcaption}
\usepackage{lmodern}
\usepackage{mdframed}
\usepackage{tcolorbox}
\usepackage{amsmath}

\lstset{
  literate=
    {é}{{\'{e}}}1 {è}{{\`{e}}}1 {ê}{{\^{e}}}1 {ë}{{\"e}}1
    {à}{{\`{a}}}1 {â}{{\^{a}}}1 {ù}{{\`{u}}}1 {û}{{\^{u}}}1
    {î}{{\^{i}}}1 {ï}{{\"i}}1   {ô}{{\^{o}}}1 {ç}{{\c{c}}}1
    {É}{{\'{E}}}1 {È}{{\`{E}}}1 {Ê}{{\^{E}}}1 {Ë}{{\"E}}1
    {À}{{\`{A}}}1 {Â}{{\^{A}}}1 {Ù}{{\`{U}}}1 {Û}{{\^{U}}}1
    {Î}{{\^{I}}}1 {Ï}{{\"I}}1   {Ô}{{\^{O}}}1 {Ç}{{\c{C}}}1
    {ë}{{\"e}}1   {î}{{\^i}}1
    {→}{{$\rightarrow$}}1 {—}{{---}}1 {–}{{--}}1
}

\newcommand{\screenshot}[2][0.85]{%
  \IfFileExists{#2}{%
    \includegraphics[width=#1\textwidth]{#2}%
  }{%
    \fbox{\begin{minipage}{#1\textwidth}%
      \centering\vspace{1cm}%
      \textit{[Capture d'\'ecran \`{a} ins\'erer]}%
      \vspace{1cm}%
    \end{minipage}}%
  }%
}

\geometry{
  top=2.5cm,
  bottom=2.5cm,
  left=2.5cm,
  right=2.5cm
}

\definecolor{codegreen}{rgb}{0.0, 0.5, 0.0}
\definecolor{codegray}{rgb}{0.5, 0.5, 0.5}
\definecolor{codeblue}{rgb}{0.13, 0.29, 0.53}
\definecolor{codepurple}{rgb}{0.55, 0.0, 0.55}
\definecolor{backcolour}{rgb}{0.97, 0.97, 0.97}
\definecolor{keywordcolor}{rgb}{0.0, 0.3, 0.6}
\definecolor{stringcolor}{rgb}{0.64, 0.08, 0.08}
\definecolor{bordercolor}{rgb}{0.7, 0.7, 0.7}

\lstdefinestyle{terraform}{
  backgroundcolor=\color{backcolour},
  commentstyle=\color{codegreen}\itshape,
  keywordstyle=\color{keywordcolor}\bfseries,
  stringstyle=\color{stringcolor},
  basicstyle=\ttfamily\footnotesize,
  breakatwhitespace=false,
  breaklines=true,
  captionpos=b,
  keepspaces=true,
  numbers=left,
  numberstyle=\tiny\color{codegray},
  numbersep=8pt,
  showspaces=false,
  showstringspaces=false,
  showtabs=false,
  tabsize=2,
  frame=single,
  rulecolor=\color{bordercolor},
  xleftmargin=15pt,
  framexleftmargin=15pt,
  morekeywords={terraform, required_providers, source, version, provider, resource,
                variable, output, default, type, description, string, number, bool,
                depends_on, ports, internal, external, env, build, context, dockerfile,
                triggers, keep_locally, name, image, sensitive, value, true, false,
                path, module, filemd5}
}

\lstdefinestyle{dockerfile}{
  backgroundcolor=\color{backcolour},
  commentstyle=\color{codegreen}\itshape,
  keywordstyle=\color{keywordcolor}\bfseries,
  stringstyle=\color{stringcolor},
  basicstyle=\ttfamily\footnotesize,
  breakatwhitespace=false,
  breaklines=true,
  captionpos=b,
  keepspaces=true,
  numbers=left,
  numberstyle=\tiny\color{codegray},
  numbersep=8pt,
  showspaces=false,
  showstringspaces=false,
  showtabs=false,
  tabsize=2,
  frame=single,
  rulecolor=\color{bordercolor},
  xleftmargin=15pt,
  framexleftmargin=15pt,
  morekeywords={FROM, RUN, COPY, ADD, EXPOSE, CMD, ENTRYPOINT, ENV, WORKDIR, ARG, LABEL}
}

\lstdefinestyle{yaml}{
  backgroundcolor=\color{backcolour},
  commentstyle=\color{codegreen}\itshape,
  keywordstyle=\color{keywordcolor}\bfseries,
  stringstyle=\color{stringcolor},
  basicstyle=\ttfamily\footnotesize,
  breakatwhitespace=false,
  breaklines=true,
  captionpos=b,
  keepspaces=true,
  numbers=left,
  numberstyle=\tiny\color{codegray},
  numbersep=8pt,
  showspaces=false,
  showstringspaces=false,
  showtabs=false,
  tabsize=2,
  frame=single,
  rulecolor=\color{bordercolor},
  xleftmargin=15pt,
  framexleftmargin=15pt,
  morekeywords={name, on, push, branches, jobs, runs-on, steps, uses, with, run, if, id,
                defaults, working-directory, continue-on-error, outcome, secrets}
}

\lstdefinestyle{bash}{
  backgroundcolor=\color{backcolour},
  commentstyle=\color{codegreen}\itshape,
  keywordstyle=\color{keywordcolor}\bfseries,
  stringstyle=\color{stringcolor},
  basicstyle=\ttfamily\footnotesize,
  breakatwhitespace=false,
  breaklines=true,
  captionpos=b,
  keepspaces=true,
  numbers=none,
  showspaces=false,
  showstringspaces=false,
  frame=single,
  rulecolor=\color{bordercolor},
  xleftmargin=5pt
}

\pagestyle{fancy}
\fancyhf{}
\fancyhead[L]{\textit{TP IaC \& CI/CD — Terraform + Docker}}
\fancyhead[R]{\textit{Génie Logiciel / DevOps}}
\fancyfoot[C]{\thepage}
\renewcommand{\headrulewidth}{0.5pt}

\hypersetup{
  colorlinks=true,
  linkcolor=codeblue,
  urlcolor=codeblue,
  pdftitle={Compte Rendu TP IaC et CI/CD},
  pdfauthor={Mootez Souilem}
}

\titleformat{\chapter}[display]
  {\normalfont\huge\bfseries\color{codeblue}}
  {\chaptertitlename\ \thechapter}{20pt}{\Huge}

\titleformat{\section}
  {\normalfont\Large\bfseries\color{codeblue}}
  {\thesection}{1em}{}

\titleformat{\subsection}
  {\normalfont\large\bfseries\color{black}}
  {\thesubsection}{1em}{}

\begin{document}

% ============================================================
% PAGE DE GARDE
% ============================================================
\begin{titlepage}
  \centering
  \vspace*{1cm}

  {\scshape\Large Travaux Pratiques --- DevOps \par}
  \vspace{0.5cm}
  {\large Génie Logiciel / Ingénierie Informatique \par}
  \vspace{2cm}

  \rule{\linewidth}{1.5pt} \\[0.4cm]
  {\Huge\bfseries Infrastructure as Code \& \\ Pipelines CI/CD \par}
  \rule{\linewidth}{1.5pt} \\[0.5cm]

  \vspace{1cm}
  {\Large\itshape Déploiement Automatisé avec Terraform et Docker \par}
  \vspace{2cm}

  \begin{tabular}{rl}
    \textbf{Étudiant :}   & Mootez Souilem \\[4pt]
    \textbf{Email :}      & mootez.souilem@bysur.com \\[4pt]
    \textbf{Date :}       & 16 avril 2026 \\[4pt]
    \textbf{Outils :}     & Docker, Terraform, GitHub Actions \\[4pt]
    \textbf{Durée :}      & 6 heures \\
  \end{tabular}

  \vfill

  \begin{tcolorbox}[colback=codeblue!10, colframe=codeblue, width=0.85\textwidth]
    \centering
    \textbf{Objectif :} Maîtriser le Cycle de Vie du Déploiement (DLC) complet via l'IaC\\
    en orchestrant des conteneurs Docker avec Terraform et en automatisant\\
    le processus à travers un pipeline CI/CD GitHub Actions.
  \end{tcolorbox}

\end{titlepage}

\tableofcontents
\listoffigures
\clearpage

% ============================================================
% INTRODUCTION
% ============================================================
\chapter*{Introduction}
\addcontentsline{toc}{chapter}{Introduction}

L'Infrastructure as Code (IaC) représente une approche fondamentale du DevOps moderne :
elle consiste à décrire et gérer l'infrastructure informatique au travers de fichiers de
configuration versionnés, lisibles par des humains et interprétables par des machines.
Contrairement à une gestion manuelle et ad hoc, l'IaC garantit la \textbf{reproductibilité},
la \textbf{traçabilité} et la \textbf{cohérence} des environnements.

Ce TP s'articule en deux grandes parties :

\begin{enumerate}
  \item \textbf{Partie I} --- Déploiement d'une infrastructure locale composée d'un conteneur
        PostgreSQL et d'une application web Nginx, orchestrés par Terraform via le provider
        Docker.
  \item \textbf{Partie II} --- Automatisation complète du Cycle de Vie du Déploiement au sein
        d'un pipeline CI/CD GitHub Actions déclenché à chaque \textit{push} sur la branche
        principale.
\end{enumerate}

\noindent L'environnement de développement utilisé est le suivant :
\begin{itemize}
  \item Système d'exploitation : Ubuntu 22.04 LTS
  \item Terraform CLI : version 1.7.x
  \item Docker Engine : version 24.x
  \item GitHub Actions : runner \texttt{ubuntu-latest}
\end{itemize}

% ============================================================
% PARTIE I
% ============================================================
\chapter{Partie I : Infrastructure as Code Locale}

\section{Structure du Projet}

Le projet suit une organisation claire et modulaire, où chaque fichier joue un rôle précis
dans la chaîne de déploiement :

\begin{mdframed}[backgroundcolor=backcolour, linecolor=bordercolor]
\begin{verbatim}
tp-iac-local/
|-- main.tf           (Provider Docker + definition des ressources)
|-- variables.tf      (Parametres configurables : ports, credentials)
|-- outputs.tf        (Informations de sortie apres deploiement)
|-- Dockerfile_app    (Blueprint de l'image Nginx/Alpine)
+-- .gitignore        (Exclusion des fichiers d'etat Terraform)
\end{verbatim}
\end{mdframed}

\section{Conteneurisation de l'Application (Dockerfile\_app)}

Le \texttt{Dockerfile\_app} construit une image légère basée sur Alpine Linux,
dans laquelle Nginx est installé pour servir une page de validation statique.

\begin{lstlisting}[style=dockerfile, caption={Dockerfile\_app --- Image Nginx/Alpine}]
FROM alpine:latest

RUN apk update && \
    apk add --no-cache nginx

RUN mkdir -p /var/www/localhost && \
    echo "<h1>Application Deployed via Terraform IaC!</h1>" > \
    /var/www/localhost/index.html

RUN mkdir -p /run/nginx

EXPOSE 80

CMD ["nginx", "-g", "daemon off;"]
\end{lstlisting}

\textbf{Choix techniques :}
\begin{itemize}
  \item \texttt{alpine:latest} : image de base minimaliste ($\approx$ 5 Mo) réduisant la surface d'attaque.
  \item \texttt{--no-cache} : évite la mise en cache du gestionnaire de paquets pour alléger l'image finale.
  \item \texttt{daemon off;} : Nginx s'exécute en premier plan, permettant à Docker de gérer le cycle de vie du processus.
  \item Création explicite de \texttt{/run/nginx} : nécessaire sur Alpine pour que Nginx puisse écrire son PID.
\end{itemize}

\section{Définition Terraform}

\subsection{variables.tf}

Ce fichier centralise tous les paramètres configurables de l'infrastructure,
rendant le code \textit{DRY} (Don't Repeat Yourself) et facilement injectable
par le pipeline CI/CD.

\begin{lstlisting}[style=terraform, caption={variables.tf --- Paramètres configurables}]
variable "db_name" {
  description = "Nom de la base de données PostgreSQL."
  type        = string
  default     = "devops_db"
}

variable "db_user" {
  description = "Nom d'utilisateur PostgreSQL."
  type        = string
  default     = "devops_user"
}

variable "db_password" {
  description = "Mot de passe PostgreSQL."
  type        = string
  default     = "strongpassword123"
  sensitive   = true
}

variable "app_port_external" {
  description = "Port externe pour accéder à l'application web."
  type        = number
  default     = 8080
}
\end{lstlisting}

L'attribut \texttt{sensitive = true} sur \texttt{db\_password} empêche Terraform
d'afficher la valeur en clair dans les logs du terminal et du pipeline CI/CD.

\subsection{main.tf}

Le fichier central déclare le provider Docker et définit l'ensemble des ressources
à provisionner.

\begin{lstlisting}[style=terraform, caption={main.tf --- Configuration principale Terraform}]
terraform {
  required_providers {
    docker = {
      source  = "kreuzwerker/docker"
      version = "~> 3.0.1"
    }
  }
}

provider "docker" {}

resource "docker_image" "postgres_image" {
  name         = "postgres:latest"
  keep_locally = true
}

resource "docker_container" "db_container" {
  name  = "tp-db-postgres"
  image = docker_image.postgres_image.image_id

  ports {
    internal = 5432
    external = 5432
  }

  env = [
    "POSTGRES_USER=${var.db_user}",
    "POSTGRES_PASSWORD=${var.db_password}",
    "POSTGRES_DB=${var.db_name}",
  ]
}

resource "docker_image" "app_image" {
  name = "tp-web-app:latest"

  build {
    context    = path.module
    dockerfile = "Dockerfile_app"
  }

  triggers = {
    dockerfile_hash = filemd5("${path.module}/Dockerfile_app")
  }
}

resource "docker_container" "app_container" {
  name  = "tp-app-web"
  image = docker_image.app_image.image_id

  depends_on = [
    docker_container.db_container
  ]

  ports {
    internal = 80
    external = var.app_port_external
  }
}
\end{lstlisting}

\textbf{Points clés du \texttt{main.tf} :}
\begin{itemize}
  \item \texttt{keep\_locally = true} : préserve l'image PostgreSQL localement après \texttt{terraform destroy} pour éviter un re-téléchargement inutile.
  \item \texttt{image\_id} : attribut correct pour le provider v3.x (l'attribut \texttt{latest} a été déprécié).
  \item \texttt{triggers} : force la reconstruction de l'image Docker si le \texttt{Dockerfile\_app} est modifié.
  \item \texttt{depends\_on} : garantit l'ordre de démarrage --- la base de données doit être opérationnelle avant le lancement de l'application.
\end{itemize}

\subsection{outputs.tf}

\begin{lstlisting}[style=terraform, caption={outputs.tf --- Sorties du déploiement}]
output "db_container_name" {
  description = "Nom du conteneur de la base de données."
  value       = docker_container.db_container.name
}

output "app_access_url" {
  description = "URL d'accès à l'application web."
  value       = "http://localhost:${docker_container.app_container.ports[0].external}"
}

output "db_container_id" {
  description = "ID du conteneur PostgreSQL."
  value       = docker_container.db_container.id
}

output "app_container_id" {
  description = "ID du conteneur de l'application web."
  value       = docker_container.app_container.id
}
\end{lstlisting}

\section{Exécution du Cycle de Vie du Déploiement (DLC)}

\subsection{Étape 1 : Initialisation (\texttt{terraform init})}

La commande \texttt{terraform init} télécharge et installe le provider Docker
depuis le registre Terraform.

\begin{lstlisting}[style=bash, caption={Commande d'initialisation}]
$ cd tp-iac-local
$ terraform init
\end{lstlisting}

\begin{figure}[H]
  \centering
  \screenshot{screenshots/terraform_init.png}
  \caption{Résultat de \texttt{terraform init} --- Téléchargement du provider \texttt{kreuzwerker/docker}}
  \label{fig:terraform-init}
\end{figure}

Le message \texttt{Terraform has been successfully initialized!} confirme que
le provider est installé dans le répertoire \texttt{.terraform/} et que le fichier
de verrou \texttt{.terraform.lock.hcl} a été généré.

\subsection{Étape 2 : Planification (\texttt{terraform plan})}

\begin{lstlisting}[style=bash, caption={Commande de planification}]
$ terraform plan
\end{lstlisting}

\begin{figure}[H]
  \centering
  \screenshot{screenshots/terraform_plan.png}
  \caption{Résultat de \texttt{terraform plan} --- Aperçu des 4 ressources à créer}
  \label{fig:terraform-plan}
\end{figure}

Terraform affiche un plan d'exécution détaillant les \textbf{4 ressources} à créer :
\begin{enumerate}
  \item \texttt{docker\_image.postgres\_image} --- Pull de \texttt{postgres:latest}
  \item \texttt{docker\_container.db\_container} --- Conteneur PostgreSQL
  \item \texttt{docker\_image.app\_image} --- Build de l'image Nginx
  \item \texttt{docker\_container.app\_container} --- Conteneur de l'application web
\end{enumerate}

Le résumé \texttt{Plan: 4 to add, 0 to change, 0 to destroy} confirme qu'aucune
infrastructure existante ne sera perturbée.

\subsection{Étape 3 : Application (\texttt{terraform apply})}

\begin{lstlisting}[style=bash, caption={Commande d'application}]
$ terraform apply -auto-approve
\end{lstlisting}

\begin{figure}[H]
  \centering
  \screenshot{screenshots/terraform_apply.png}
  \caption{Résultat de \texttt{terraform apply} --- Création des 4 ressources et affichage des outputs}
  \label{fig:terraform-apply}
\end{figure}

À l'issue de l'apply, Terraform affiche les \textit{outputs} définis dans
\texttt{outputs.tf}, notamment l'URL d'accès :

\begin{mdframed}[backgroundcolor=backcolour, linecolor=bordercolor]
\begin{verbatim}
Apply complete! Resources: 4 added, 0 changed, 0 destroyed.

Outputs:
app_access_url    = "http://localhost:8080"
db_container_name = "tp-db-postgres"
\end{verbatim}
\end{mdframed}

\subsection{Étape 4 : Validation}

Après l'application, la vérification s'effectue à deux niveaux :

\subsubsection{Vérification via navigateur}

\begin{figure}[H]
  \centering
  \screenshot{screenshots/app_browser.png}
  \caption{Accès à \texttt{http://localhost:8080} --- Page de confirmation du déploiement IaC}
  \label{fig:app-browser}
\end{figure}

\subsubsection{Vérification via Docker CLI}

\begin{lstlisting}[style=bash, caption={Vérification de l'état des conteneurs}]
$ docker ps
\end{lstlisting}

\begin{figure}[H]
  \centering
  \screenshot{screenshots/docker_ps.png}
  \caption{Résultat de \texttt{docker ps} --- Deux conteneurs en état \texttt{Up}}
  \label{fig:docker-ps}
\end{figure}

Les deux conteneurs \texttt{tp-db-postgres} et \texttt{tp-app-web} apparaissent
avec le statut \texttt{Up}, confirmant le bon déroulement du déploiement.

\subsection{Étape 5 : Destruction (\texttt{terraform destroy})}

\begin{lstlisting}[style=bash, caption={Commande de destruction}]
$ terraform destroy -auto-approve
\end{lstlisting}

\begin{figure}[H]
  \centering
  \screenshot{screenshots/terraform_destroy.png}
  \caption{Résultat de \texttt{terraform destroy} --- Suppression des 4 ressources}
  \label{fig:terraform-destroy}
\end{figure}

Le message final \texttt{Destroy complete! Resources: 4 destroyed.} confirme
la suppression propre de l'ensemble des ressources provisionnées.

% ============================================================
% PARTIE II
% ============================================================
\chapter{Partie II : Pipeline CI/CD GitHub Actions}

\section{Contexte et Architecture du Pipeline}

L'objectif de cette partie est d'automatiser le Cycle de Vie du Déploiement en
intégrant les commandes Terraform dans un pipeline GitHub Actions. Ce pipeline
se déclenche automatiquement à chaque \textit{push} sur la branche \texttt{main}.


\section{Configuration du Pipeline (\texttt{main.yml})}

\begin{lstlisting}[style=yaml, caption={.github/workflows/main.yml --- Pipeline GitHub Actions}]
name: Terraform IaC Pipeline

on:
  push:
    branches:
      - main

jobs:
  terraform:
    name: Deploy Infrastructure
    runs-on: ubuntu-latest
    defaults:
      run:
        working-directory: tp-iac-local

    steps:
      - name: Checkout Repository
        uses: actions/checkout@v4

      - name: Setup Terraform
        uses: hashicorp/setup-terraform@v3
        with:
          terraform_version: "1.7.0"

      - name: Terraform Init
        run: terraform init

      - name: Terraform Plan
        id: plan
        run: terraform plan -var="db_password=${{ secrets.DB_PASSWORD }}" -out=tfplan
        continue-on-error: false

      - name: Terraform Apply
        if: steps.plan.outcome == 'success'
        run: terraform apply -auto-approve tfplan
\end{lstlisting}

\textbf{Points de conception importants :}
\begin{itemize}
  \item \texttt{working-directory: tp-iac-local} : les commandes Terraform s'exécutent
        dans le bon répertoire sans nécessiter de \texttt{cd} explicite à chaque étape.
  \item \texttt{continue-on-error: false} sur l'étape Plan : tout échec du plan
        arrête immédiatement le pipeline, empêchant un apply sur une base défaillante.
  \item \texttt{if: steps.plan.outcome == 'success'} : la condition de garde (\textit{gate})
        garantit que l'étape Apply ne s'exécute que si et seulement si le Plan a réussi.
  \item \texttt{secrets.DB\_PASSWORD} : le mot de passe de la base de données est injecté
        via les secrets GitHub, sans jamais apparaître en clair dans le code source.
  \item \texttt{-out=tfplan} : le plan est sauvegardé dans un fichier binaire, que l'étape
        Apply consomme directement --- garantissant que l'état appliqué correspond
        exactement à ce qui a été planifié.
\end{itemize}

\section{Configuration du Secret GitHub}

Avant de pousser le projet, le secret \texttt{DB\_PASSWORD} doit être créé dans les
paramètres du dépôt GitHub : \texttt{Settings > Secrets and variables > Actions > New repository secret}.

\begin{figure}[H]
  \centering
  \screenshot{screenshots/github_secret.png}
  \caption{Configuration du secret \texttt{DB\_PASSWORD} dans les paramètres GitHub Actions}
  \label{fig:github-secret}
\end{figure}

\section{Exécution et Validation du Pipeline}

\subsection{Premier déclenchement}

\begin{lstlisting}[style=bash, caption={Push initial vers la branche main}]
$ git init
$ git add .
$ git commit -m "feat: initial IaC setup with Terraform and Docker"
$ git remote add origin https://github.com/<username>/tp-iac-local.git
$ git push -u origin main
\end{lstlisting}

\begin{figure}[H]
  \centering
  \screenshot{screenshots/pipeline_run1.png}
  \caption{Exécution réussie du pipeline lors du premier push --- Toutes les étapes en vert}
  \label{fig:pipeline-run1}
\end{figure}

\subsection{Validation par modification d'une variable}

Conformément aux consignes, le mot de passe de la base de données a été modifié
dans \texttt{variables.tf} pour déclencher un second cycle de déploiement.

\begin{lstlisting}[style=bash, caption={Modification et push pour valider la mise à jour}]
$ git add variables.tf
$ git commit -m "chore: update db_password default value"
$ git push origin main
\end{lstlisting}

\begin{figure}[H]
  \centering
  \screenshot{screenshots/pipeline_run2_plan.png}
  \caption{Étape Plan du second déclenchement --- Détection du changement de variable}
  \label{fig:pipeline-run2-plan}
\end{figure}

\begin{figure}[H]
  \centering
  \screenshot{screenshots/pipeline_run2_apply.png}
  \caption{Étape Apply du second déclenchement --- Mise à jour de la ressource détectée}
  \label{fig:pipeline-run2-apply}
\end{figure}

Le pipeline détecte le changement de variable dans l'étape Plan (\texttt{1 to change})
et applique la modification lors de l'étape Apply, démontrant la boucle de réconciliation
automatique de Terraform.

% ============================================================
% QUESTIONS D'APPROFONDISSEMENT
% ============================================================
\chapter{Questions d'Approfondissement}

\section{Impact de \texttt{terraform destroy} sur \texttt{terraform.tfstate}}

\texttt{terraform.tfstate} est le \textbf{fichier d'état} de Terraform : il contient une
représentation JSON de l'infrastructure réelle telle que Terraform la connaît, incluant
les identifiants des ressources créées, leurs attributs actuels, et les dépendances
entre elles.

Lors de l'exécution de \texttt{terraform destroy} :

\begin{enumerate}
  \item Terraform lit le \texttt{terraform.tfstate} pour identifier toutes les ressources
        à supprimer et leur ordre de destruction (inversé par rapport à la création,
        en respectant les dépendances).
  \item Il contacte le provider (Docker dans notre cas) pour effectivement supprimer
        les ressources.
  \item Il met à jour le \texttt{terraform.tfstate} : les ressources supprimées sont
        retirées, laissant un fichier d'état vide (ou ne contenant que les ressources
        non détruites).
\end{enumerate}

\noindent\textbf{Rôle dans la réconciliation :} Le fichier d'état est le pont entre
\textit{l'infrastructure déclarée dans le code} et \textit{l'infrastructure réelle}.
À chaque \texttt{plan} ou \texttt{apply}, Terraform compare l'état désiré (les
fichiers \texttt{.tf}) à l'état connu (le \texttt{.tfstate}) et à l'état réel (via des
appels API au provider) pour calculer le diff minimal à appliquer. Sans ce fichier,
Terraform ne saurait pas quelles ressources il a déjà créées et tenterait de tout
recréer à chaque exécution.

\begin{tcolorbox}[colback=orange!5, colframe=orange!80!black, title=Attention en production]
  En production, le fichier \texttt{terraform.tfstate} doit être stocké dans un
  \textbf{backend distant} (S3, Azure Blob, Terraform Cloud) avec verrouillage
  (\textit{state locking}) pour éviter les corruptions en cas d'accès concurrent.
  Ne jamais committer ce fichier dans Git.
\end{tcolorbox}

\section{Immuabilité de l'Infrastructure}

L'immuabilité de l'infrastructure (\textit{Immutable Infrastructure}) est un principe
selon lequel \textbf{un serveur ou un conteneur déployé ne doit jamais être modifié
après sa création}. Toute modification passe par la création d'une nouvelle instance
qui remplace l'ancienne, au lieu de muter l'instance existante.

\begin{center}
\begin{tabular}{>{\bfseries}p{5cm} p{8cm}}
  \toprule
  Approche Mutable & Approche Immuable \\
  \midrule
  On se connecte au serveur en SSH et on met à jour le paquet Nginx &
  On reconstruit une nouvelle image Docker avec la nouvelle version, on déploie
  le nouveau conteneur, on supprime l'ancien \\
  \midrule
  Risque de \textit{configuration drift} &
  Infrastructure toujours dans un état connu et reproductible \\
  \midrule
  Retour arrière difficile &
  Retour arrière = redéployer l'image précédente \\
  \bottomrule
\end{tabular}
\end{center}

\vspace{0.5em}
\noindent Dans notre TP, l'utilisation du mécanisme \texttt{triggers} dans la ressource
\texttt{docker\_image.app\_image} illustre ce principe : toute modification du
\texttt{Dockerfile\_app} provoque la \textbf{reconstruction complète} de l'image
et le \textbf{remplacement} du conteneur, jamais sa modification en place.

\section{Rôle de l'Étape Plan comme Gate du DLC}

L'étape \texttt{terraform plan} joue le rôle de \textbf{porte de validation}
(\textit{quality gate}) essentielle dans le DLC pour plusieurs raisons :

\begin{enumerate}
  \item \textbf{Prévisibilité sans risque} : le plan simule l'intégralité des changements
        sans toucher à l'infrastructure réelle. Un opérateur peut vérifier exactement
        ce qui sera modifié, créé ou supprimé.

  \item \textbf{Détection précoce d'erreurs} : les erreurs de syntaxe, de références
        incorrectes ou de conflits de ressources sont détectées avant toute modification
        de production.

  \item \textbf{Séparation des responsabilités} : dans un contexte d'équipe, le plan
        peut être soumis à une revue humaine (\textit{plan review}) ou à une approbation
        manuelle avant l'apply --- c'est le principe du \textit{four-eyes principle}.

  \item \textbf{Intégrité du pipeline} : dans notre fichier \texttt{main.yml}, la condition
        \texttt{if: steps.plan.outcome == 'success'} garantit que l'étape Apply
        n'est jamais exécutée sur une base de plan défaillante.

  \item \textbf{Auditabilité} : le fichier \texttt{tfplan} généré constitue une preuve
        immuable de ce qui était prévu d'être appliqué, utilisable à des fins d'audit.
\end{enumerate}

\section{Réseau Docker comme Ressource Terraform}

Pour créer un réseau Docker dédié et y connecter les deux conteneurs, la ressource
Terraform à utiliser est \textbf{\texttt{docker\_network}} du provider
\texttt{kreuzwerker/docker}. Voici comment l'intégrer dans \texttt{main.tf} :

\begin{lstlisting}[style=terraform, caption={Ajout d'un réseau Docker partagé}]
resource "docker_network" "app_network" {
  name = "tp-app-network"
}

resource "docker_container" "db_container" {
  name  = "tp-db-postgres"
  image = docker_image.postgres_image.image_id

  networks_advanced {
    name = docker_network.app_network.name
  }

  env = [
    "POSTGRES_USER=${var.db_user}",
    "POSTGRES_PASSWORD=${var.db_password}",
    "POSTGRES_DB=${var.db_name}",
  ]
}

resource "docker_container" "app_container" {
  name  = "tp-app-web"
  image = docker_image.app_image.image_id

  depends_on = [docker_container.db_container]

  networks_advanced {
    name = docker_network.app_network.name
  }

  ports {
    internal = 80
    external = var.app_port_external
  }
}
\end{lstlisting}

\noindent\textbf{Avantages de cette approche :}
\begin{itemize}
  \item Les conteneurs communiquent par leur nom DNS (ex : le conteneur \texttt{app}
        peut joindre le conteneur \texttt{db} via l'hôte \texttt{tp-db-postgres}).
  \item Le port 5432 de PostgreSQL n'a plus besoin d'être exposé sur l'hôte,
        améliorant la sécurité (\textit{principle of least exposure}).
  \item L'infrastructure est plus proche d'un environnement de production réel.
\end{itemize}

% ============================================================
% CONCLUSION
% ============================================================
\chapter*{Conclusion}
\addcontentsline{toc}{chapter}{Conclusion}

Ce TP a permis de mettre en pratique l'ensemble du cycle de vie d'un déploiement
infrastructure moderne, depuis la définition déclarative des ressources avec Terraform
jusqu'à leur automatisation complète via un pipeline CI/CD GitHub Actions.

\vspace{0.5em}
\noindent Les compétences acquises couvrent trois dimensions complémentaires :

\begin{itemize}
  \item \textbf{Technique} : maîtrise du provider Docker Terraform (v3.x), gestion des
        dépendances entre ressources, build d'images Docker depuis Terraform, gestion
        des outputs et des variables sensibles.

  \item \textbf{Processus} : compréhension du Cycle de Vie du Déploiement à travers
        ses cinq étapes (Init → Plan → Apply → Validate → Destroy), avec la distinction
        cruciale entre phase de planification (prévisualisation sans risque) et phase
        d'application (mutation réelle).

  \item \textbf{Culture DevOps} : appréhension des principes fondateurs --- immuabilité
        de l'infrastructure, gestion de l'état, séparation entre code d'infrastructure
        et secrets, et automatisation comme garant de la cohérence.
\end{itemize}

\vspace{0.5em}
\noindent Ce socle constitue une base solide pour aborder des environnements cloud
réels (AWS, Azure, GCP) où Terraform est l'outil de facto pour l'orchestration
d'infrastructure à grande échelle.

\end{document}
